\Conclusion % заключение к отчёту

Выполнение выпускной квалификационной работы базировалось на анализе предметной области, сравнительном исследовании существующих решений, проектировании архитектуры системы и практической реализации программного продукта.

В ходе выполнения работы были решены следующие задачи:
\begin{compactlist}
  \item выполнен анализ предметной области планирования питания и выявлены ключевые проблемы целевой аудитории;
  \item проведен анализ аналогов и конкурентов, определены ограничения существующих подходов;
  \item спроектированы бизнес-процессы и формализованы сценарии взаимодействия пользователей с системой;
  \item сформированы бизнес-, функциональные и нефункциональные требования к разрабатываемой системе;
  \item спроектирована архитектура веб-приложения и определена ролевая модель доступа;
  \item разработана структура базы данных и описаны связи доменных сущностей;
  \item разработаны макеты и реализован пользовательский интерфейс клиентской части;
  \item реализованы серверная и клиентская части приложения, включая контуры персонализации, модерации, отзывов, уведомлений и ИИ-ассистента;
  \item выполнено тестирование ключевых пользовательских и административных сценариев;
  \item описаны и реализованы подходы к контейнеризации и развертыванию системы.
\end{compactlist}

По результатам выполненной работы разработанное веб-приложение FoodPlanner демонстрирует работоспособность заявленных сценариев и может использоваться в качестве основы для дальнейшего внедрения и развития.

В качестве перспектив развития рассматриваются расширение аналитического контура, развитие инструментов модерации и администрирования, а также усиление персонализации рекомендаций на основе пользовательского поведения.

Таким образом, поставленные задачи решены, а цель выпускной квалификационной работы --- разработка веб-приложения для планирования питания с учетом предпочтений пользователей --- достигнута.

